% from ega4-2-fr.tex
\oldpage[IV-2]{112}

la définition du foncteur $\Gamma$ et du fait qu'il est exact à gauche.
Comme l'homomorphisme $(\rho_{X/Z})_V$ est le composé
\begin{equation}
\Gamma(V,\mathcal F)\longrightarrow
\Gamma(V\cap(X-T),\mathcal F)\longrightarrow
\varinjlim_\alpha\Gamma(V\cap U_\alpha,\mathcal F)
\tag{5.9.8.1}\label{IV.2.s5.9.8.1-fr}
\end{equation}
pour toute partie fermée $T\subset Z$, si $(\rho_{X/Z})_V$ est injectif, il
en est de même de
$\Gamma(V,\mathcal F)\to\Gamma(V\cap(X-T),\mathcal F)$; d'autre part, le
fait que b') implique b) résulte de la définition d'une limite inductive. Il
reste à montrer que si $\rho_{X/Z}$ est bijectif, il en est de même de
$\Gamma(V,\mathcal F)\to\Gamma(V\cap(X-T),\mathcal F)$, et pour cela il
suffit, en vertu de (5.9.8.1), de voir que si $U'\subset U$ sont deux
ouverts contenus dans $V$ et contenant $V\cap Z$, l'homomorphisme de
restriction $\Gamma(U,\mathcal F)\to\Gamma(U',\mathcal F)$ est injectif;
mais cela résulte de ce que $\rho_{X/Z}$ est injectif, en remplaçant dans ce
qui précède $V$ par $U$ et $V\cap(X-T)$ par $U'$.

\begin{definition}[5.9.9]
\label{IV.2.s5.9.9-fr}
Sous les hypothèses de (5.9.8), on dit que $\mathcal F$ est $Z$-pur
(resp. $Z$-clos) si l'homomorphisme $\rho_{X/Z}$ est injectif
(resp. bijectif).
\end{definition}

Si $X=\operatorname{Spec}(A)$ est affine, $\mathcal F=\widetilde M$, où
$M$ est un $A$-module, on dira que $M$ est $Z$-pur (resp. $Z$-clos) lorsque
$\mathcal F$ est $Z$-pur (resp. $Z$-clos).

On dit que $\mathcal F$ est $Z$-pur (resp. $Z$-clos) \emph{en un point}
$x\in X$ si (avec les notations de (5.9.6)) $\widetilde{\mathcal F}_x$ est
$Z_x$-pur (resp. $Z_x$-clos); il revient au même, en vertu de (5.9.6), de
dire que l'homomorphisme canonique
\[
\mathcal F_x\longrightarrow(\mathcal H^0_{X/Z}(\mathcal F))_x
\]
est injectif (resp. bijectif).

On notera que pour tout $x\in X-Z$, $\mathcal F$ est $Z$-clos au point $x$,
en vertu de (5.9.8).

\begin{corollary}[5.9.10]
\label{IV.2.s5.9.10-fr}
\begin{enumerate}
\item[(i)] Soit $(V_\lambda)$ un recouvrement ouvert de $X$. Pour que
$\mathcal F$ soit $Z$-pur (resp. $Z$-clos), il faut et il suffit que pour
tout $\lambda$, $\mathcal F|V_\lambda$ soit $(Z\cap V_\lambda)$-pur
(resp. $(Z\cap V_\lambda)$-clos).
\item[(ii)] Soit $Z'$ une partie de $Z$ stable par spécialisation. Si
$\mathcal F$ est $Z$-pur (resp. $Z$-clos), il est $Z'$-pur
(resp. $Z'$-clos).
\end{enumerate}
\end{corollary}

Cela résulte aussitôt de (5.9.8, b')).

\begin{proposition}[5.9.11]
\label{IV.2.s5.9.11-fr}
Sous les hypothèses de (5.9.8), les $\mathcal O_X$-Modules
$\operatorname{Ker}(\rho_{X/Z})$ et $\operatorname{Coker}(\rho_{X/Z})$ ont
leur support contenu dans $Z$, et le $\mathcal O_X$-Module
$\mathcal H^0_{X/Z}(\mathcal F)$ est $Z$-clos. En outre, si
$u:\mathcal F\to\mathcal F'$ est un homomorphisme de
$\mathcal O_X$-Modules tel que $\mathcal F'$ soit $Z$-clos, $u$ se factorise
d'une seule manière en
\[
\mathcal F\xrightarrow{\rho_{X/Z}}\mathcal H^0_{X/Z}(\mathcal F)
\xrightarrow{v}\mathcal F'.
\]
Si en outre les supports de $\operatorname{Ker}(u)$ et de
$\operatorname{Coker}(u)$ sont contenus dans $Z$, $v$ est un isomorphisme.
\end{proposition}

La première assertion signifie que, pour tout $x\in X-Z$, on a
\[
\operatorname{Ker}(\rho_{X/Z})_x
=\operatorname{Coker}(\rho_{X/Z})_x=0,
\]
i.e. que $(\rho_{X/Z})_x$ est bijectif, ou encore que $\mathcal F$ est
$Z$-pur en $x$, comme il a été signalé plus haut (5.9.9).

Pour montrer que $\mathcal H^0_{X/Z}(\mathcal F)$ est $Z$-clos, il s'agit de
voir que pour un ouvert noethérien $V$ de $X$,
\[
\varinjlim_\beta
\Gamma(V\cap U_\beta,\mathcal H^0_{X/Z}(\mathcal F))
\]
est égale à $\Gamma(V,\mathcal H^0_{X/Z}(\mathcal F))$; mais par définition
\[
\Gamma(V\cap U_\beta,\mathcal H^0_{X/Z}(\mathcal F))
=\varinjlim_\alpha\Gamma(V\cap U_\alpha\cap U_\beta,\mathcal F).
\]
Or, la famille double $(U_\alpha\cap U_\beta)$ est filtrante décroissante,
et notre assertion résulte de (5.9.1) et du théorème de la double limite
inductive.

